\section{Evaluation}

\subsection{Impact on Score Improvement}

The primary evaluation question: does the P1/P2/P3 classification direct revision effort effectively?

\begin{table}[h]
\centering
\begin{tabular}{lccc}
\toprule
Priority & Score Impact & \% of Total Improvement & Effort (avg days) \\
\midrule
P1 addressed & +1.3 points & 72\% & 5--7 \\
P2 addressed & +0.4 points & 22\% & 3--5 \\
P3 addressed & +0.1 points & 6\% & 1--2 \\
\bottomrule
\end{tabular}
\caption{Score impact by priority level.}
\end{table}

P1 items account for nearly three-quarters of achievable improvement despite comprising only 23\% of total issues. This validates the triage: the classification successfully identifies the issues that matter most.

\subsection{Deduplication Accuracy}

We evaluate deduplication by manually annotating 100 issue pairs from 10 synthesis cycles:

\begin{itemize}
    \item \textbf{Precision}: 91\% of merged issues are genuine duplicates
    \item \textbf{Recall}: 78\% of actual duplicates are identified
    \item \textbf{False merges}: 9\% of merges combine distinct issues (Type I error)
    \item \textbf{Missed merges}: 22\% of duplicates remain separate (Type II error)
\end{itemize}

The system errs toward under-merging (keeping separate) rather than over-merging, which is the safer direction: duplicate issues waste some author time on re-reading, while falsely merged issues might cause a concern to be overlooked.

\subsection{Synthesis Quality}

Comparing synthesis documents against the underlying reviews:
\begin{itemize}
    \item \textbf{Coverage}: 94\% of major issues in reviews appear in synthesis P1 or P2
    \item \textbf{Attribution accuracy}: 98\% of reviewer attributions are correct
    \item \textbf{Priority agreement}: Manual classification by the author agrees with automated classification on 87\% of items
\end{itemize}
